% Options for packages loaded elsewhere
\PassOptionsToPackage{unicode}{hyperref}
\PassOptionsToPackage{hyphens}{url}
\documentclass[
  ignorenonframetext,
]{beamer}
\newif\ifbibliography
\usepackage{pgfpages}
\setbeamertemplate{caption}[numbered]
\setbeamertemplate{caption label separator}{: }
\setbeamercolor{caption name}{fg=normal text.fg}
\beamertemplatenavigationsymbolsempty
% remove section numbering
\setbeamertemplate{part page}{
  \centering
  \begin{beamercolorbox}[sep=16pt,center]{part title}
    \usebeamerfont{part title}\insertpart\par
  \end{beamercolorbox}
}
\setbeamertemplate{section page}{
  \centering
  \begin{beamercolorbox}[sep=12pt,center]{section title}
    \usebeamerfont{section title}\insertsection\par
  \end{beamercolorbox}
}
\setbeamertemplate{subsection page}{
  \centering
  \begin{beamercolorbox}[sep=8pt,center]{subsection title}
    \usebeamerfont{subsection title}\insertsubsection\par
  \end{beamercolorbox}
}
% Prevent slide breaks in the middle of a paragraph
\widowpenalties 1 10000
\raggedbottom
\AtBeginPart{
  \frame{\partpage}
}
\AtBeginSection{
  \ifbibliography
  \else
    \frame{\sectionpage}
  \fi
}
\AtBeginSubsection{
  \frame{\subsectionpage}
}
\usepackage{iftex}
\ifPDFTeX
  \usepackage[T1]{fontenc}
  \usepackage[utf8]{inputenc}
  \usepackage{textcomp} % provide euro and other symbols
\else % if luatex or xetex
  \usepackage{unicode-math} % this also loads fontspec
  \defaultfontfeatures{Scale=MatchLowercase}
  \defaultfontfeatures[\rmfamily]{Ligatures=TeX,Scale=1}
\fi
\usepackage{lmodern}
\ifPDFTeX\else
  % xetex/luatex font selection
\fi
% Use upquote if available, for straight quotes in verbatim environments
\IfFileExists{upquote.sty}{\usepackage{upquote}}{}
\IfFileExists{microtype.sty}{% use microtype if available
  \usepackage[]{microtype}
  \UseMicrotypeSet[protrusion]{basicmath} % disable protrusion for tt fonts
}{}
\makeatletter
\@ifundefined{KOMAClassName}{% if non-KOMA class
  \IfFileExists{parskip.sty}{%
    \usepackage{parskip}
  }{% else
    \setlength{\parindent}{0pt}
    \setlength{\parskip}{6pt plus 2pt minus 1pt}}
}{% if KOMA class
  \KOMAoptions{parskip=half}}
\makeatother
\setlength{\emergencystretch}{3em} % prevent overfull lines
\providecommand{\tightlist}{%
  \setlength{\itemsep}{0pt}\setlength{\parskip}{0pt}}
\usepackage{bookmark}
\IfFileExists{xurl.sty}{\usepackage{xurl}}{} % add URL line breaks if available
\urlstyle{same}
\hypersetup{
  pdftitle={Day 1 Recap},
  hidelinks,
  pdfcreator={LaTeX via pandoc}}

\title{Day 1 Recap}
\author{}
\date{\vspace{-2.5em}}

\begin{document}
\frame{\titlepage}

\begin{frame}{Workflow}
\phantomsection\label{workflow}
\begin{center}\includegraphics[width=0.6\linewidth]{../Images/workflow2} \end{center}
\end{frame}

\begin{frame}{Data set}
\phantomsection\label{data-set}
\begin{itemize}
\tightlist
\item
  Data set:
  \href{https://www.nature.com/articles/s41598-020-64929-x}{CaronBourque2020}:
  Data from Childhood acute lymphoblastic leukemia (cALL)
\item
  cells: Bone Marrow Mononuclear cells (BMMCs)

  \begin{itemize}
  \tightlist
  \item
    12 samples
  \item
    4 Sample groups

    \begin{itemize}
    \tightlist
    \item
      HHD: The high hyper diploid cases (51--67 chromosomes).

      \begin{itemize}
      \tightlist
      \item
        Two replicates.
      \end{itemize}
    \item
      PBMMC: healthy pediatric BMMC.

      \begin{itemize}
      \tightlist
      \item
        Four replicates.
      \item
        There are two PBMMC\_1 samples. These are two libraries from the
        same sample material.
      \end{itemize}
    \item
      ETV6-RUNX1: ETV6/RUNX1 rearrangement

      \begin{itemize}
      \tightlist
      \item
        Four replicates
      \end{itemize}
    \item
      Pre-T: Pre-T ALL

      \begin{itemize}
      \tightlist
      \item
        Two replicates
      \end{itemize}
    \end{itemize}
  \end{itemize}
\item
  Aim: characterize the heterogeneity of gene expression at the cell
  level, within and between patients
\end{itemize}
\end{frame}

\begin{frame}{10x library file structure}
\phantomsection\label{x-library-file-structure}
The 10x library contains four pieces of information, in the form of DNA
sequences, for each ``read''.

\begin{itemize}
\tightlist
\item
  \textbf{sample index} - identifies the library, with one or two
  indexes per sample
\item
  \textbf{10x barcode} - identifies the droplet in the library
\item
  \textbf{UMI} - identifies the transcript molecule within a cell and
  gene
\item
  \textbf{insert} - the transcript molecule
\end{itemize}

\begin{center}\includegraphics[width=0.6\linewidth]{../Images/tenxLibStructureV3dual2} \end{center}
\end{frame}

\begin{frame}{Cell Ranger}
\phantomsection\label{cell-ranger}
\begin{itemize}
\tightlist
\item
  10x Cell Ranger - This not only carries out the alignment and feature
  counting, but will also:

  \begin{itemize}
  \tightlist
  \item
    Call cells
  \item
    Generates counts matrix
  \item
    Generate a summary report in html format
  \item
    Generate a ``cloupe'' file
  \end{itemize}
\end{itemize}
\end{frame}

\begin{frame}{Cell Ranger references}
\phantomsection\label{cell-ranger-references}
cellranger mkref\\
--fasta=\{GENOME FASTA\}\\
--genes=\{ANNOTATION GTF\}\\
--genome=\{OUTPUT FOLDER FOR INDEX\}\\
--nthreads=\{CPUS\}
\end{frame}

\begin{frame}{Running cellranger count}
\phantomsection\label{running-cellranger-count}
cellranger count --id=\{OUTPUT\_SAMPLE\_NAME\}\\
--transcriptome=\{DIRECTORY\_WITH\_REFERENCE\}\\
--fastqs=\{DIRECTORY\_WITH\_FASTQ\_FILES\}\\
--sample=\{NAME\_OF\_SAMPLE\_IN\_FASTQ\_FILES\}\\
--localcores=\{NUMBER\_OF\_CPUS\}\\
--localmem=\{RAM\_MEMORY\}
\end{frame}

\begin{frame}[fragile]{Cell Ranger outputs}
\phantomsection\label{cell-ranger-outputs}
The contents of the \texttt{outs} directory are:

\begin{center}\includegraphics[width=0.75\linewidth]{../Images/CellRangerOutputOuts} \end{center}
\end{frame}

\begin{frame}{Not every droplet is useble}
\phantomsection\label{not-every-droplet-is-useble}
\begin{center}\includegraphics[width=1\linewidth,height=0.8\textheight]{../Images/droplet_overview} \end{center}
\end{frame}

\begin{frame}[fragile]{Quality Control overview}
\phantomsection\label{quality-control-overview}
\begin{itemize}
\tightlist
\item
  Aim of QC is \ldots{}

  \begin{itemize}
  \tightlist
  \item
    To remove undetected genes
  \item
    To remove empty droplets
  \item
    To remove droplets with dead cells
  \item
    To remove Doublet/multiplet
  \item
    Ultimately To filter the data to only include true cells that are of
    high quality
  \end{itemize}
\end{itemize}

\begin{frame}{The \emph{Seurat} object}
\phantomsection\label{the-seurat-object}
\end{frame}

\begin{frame}[fragile]{Useful commands in QC}
\phantomsection\label{useful-commands-in-qc}
\begin{itemize}
\item
  access counts: \texttt{seurat\_object@assays\$RNA\$counts}
\item
  access cell metadata: \texttt{seurat\_object@meta.data}
\end{itemize}
\end{frame}

\begin{frame}{QC parameters}
\phantomsection\label{qc-parameters}
\begin{itemize}
\tightlist
\item
  The library size
\item
  Number of expressed genes in each cell
\item
  proportion of UMIs mapped to genes in the mitochondrial genome
\end{itemize}
\end{frame}
\end{frame}

\end{document}
